\documentclass[11pt,a4paper]{extarticle}

\usepackage{kotex}
\usepackage[margin=18mm]{geometry}
\usepackage{amsmath,amssymb,mathtools}
\usepackage{array,booktabs,longtable,tabularx}
\usepackage[table]{xcolor}
\usepackage{enumitem}
\usepackage[hidelinks]{hyperref}
\usepackage{xurl}

\setlength{\parindent}{0pt}
\setlength{\parskip}{5pt}
\setlist[itemize]{leftmargin=18pt,itemsep=2pt,topsep=2pt}
\setlist[enumerate]{leftmargin=22pt,itemsep=3pt,topsep=3pt}
\setcounter{tocdepth}{1}

\newcommand{\CM}{\mathrm{CM}}
\newcommand{\Hash}{\mathrm{H}}
\newcommand{\MT}{\mathrm{MT}}
\newcommand{\rvMT}{\mathrm{rvMT}}
\newcommand{\State}{\mathbf{s}}

\title{zkDPP POC 시나리오 및 성능 평가 보고서}
\author{Canonical Scenario와 34개 Benchmark Profile 해설}
\date{2026-07-23}
\renewcommand{\contentsname}{목차}

\begin{document}
\maketitle

\section*{보고서 목적}

이 문서는 다음 원시 artifact를 사람이 읽고 해석할 수 있는 형태로 재구성한다.

\begin{itemize}
  \item \path{testdata/scenario/canonical.json}: POC에서 실행한 공급망 시나리오
  \item \path{benchmarks/summary.csv}: State 길이와 Process 입출력 크기에 따른 circuit 성능
  \item \path{benchmarks/anvil-transaction-gas.json}: 실제 Anvil transaction gas와 calldata
  \item \path{benchmarks/anvil-e2e-time.json}: live proof를 포함한 5회 E2E 시간
\end{itemize}

원시 JSON의 commitment, nullifier와 Merkle root 숫자를 다시 나열하는 것이 목적이 아니다.
대신 ``어떤 물품이 어떤 Event를 거쳤는가'', ``각 proof가 무엇을 검증했는가'',
``측정값이 어느 실행 계층의 비용인가''를 설명한다.

\tableofcontents
\newpage

\section{POC 한눈에 보기}

\begin{center}
  \begin{tabularx}{\textwidth}{@{}p{0.22\textwidth}X@{}}
    \toprule
    항목                & POC 설정                                                               \\
    \midrule
    Proof system      & PLONK-KZG on BLS12-381                                               \\
    Circuit framework & Go, gnark v0.15.0, gnark-crypto v0.20.1                              \\
    Hash              & BLS12-381 scalar field용 Poseidon2                                    \\
    Tree              & commitment용 $\MT$와 voucher용 $\rvMT$, 모두 depth 32                     \\
    StateVector       & $(total\_kg, carbon\_kgCO2e, recycled\_kg)$, 기본 길이 3                 \\
    숫자 표현             & Quantity는 integer, State는 $10^9$을 곱한 nano-scale unsigned integer     \\
    구현 Event          & Entry, Transfer, Proceed, Recall, Merge, Split, Process, Exit        \\
    Process           & Foil 1-to-2 두 번과 Cell 3-to-3 한 번; 성능 profile은 1-to-1, 2-to-2, 3-to-3 \\
    Contract 환경       & Solidity 0.8.30, Foundry v1.7.1, Prague EVM                          \\
    Transaction 환경    & Docker Anvil, chain ID 31337, block gas limit 30,000,000             \\
    POC SRS           & \texttt{unsafekzg.NewSRS}; production setup으로 사용하지 않음                \\
    \bottomrule
  \end{tabularx}
\end{center}

\subsection*{프로토콜 개요}

\begin{center}
  \begin{tabular}{ll}
    \hline
    Event    & Protocol                                        \\
    \hline
    Entry    & $\varnothing\rightarrow cm_{new}$                         \\
    Transfer & $cm_{old}\rightarrow (cm_{change},rv_{new})$              \\
    Proceed  & $rv_{old}\rightarrow cm_{new}$                            \\
    Recall   & $rv_{old}\rightarrow cm_{new}$                            \\
    Merge    & $(cm_{old,1},cm_{old,2})\rightarrow cm_{new}$             \\
    Split    & $cm_{old}\rightarrow (cm_{new,1},cm_{new,2})$             \\
    Process  & $(cm_{old,i})_{i=1}^{m}\rightarrow(cm_{new,j})_{j=1}^{n}$ \\
    Exit     & $cm_{old}\rightarrow\varnothing$                          \\
    \hline
  \end{tabular}
\end{center}

\[
  \begin{aligned}
    \mathit{DocumentInfo}
     & = (ProductName,LotId,Unit)                         \\[5pt]
    cm
     & = \Hash(\mathit{DocumentHash},q,\mathbf{s},addr,o) \\
    \mathbf{s}
     & = (total_{kg},carbon_{kgCO2e},recycled_{kg})       \\[5pt]
    nf
     & = \Hash(sk_{own},cm)                               \\[5pt]
    rv
     & = \Hash(
    \mathit{DocumentHash},q,\mathbf{s},
    addr_{sender},addr_{receiver},\Delta_{epoch},o_{rv}
    )                                                     \\[5pt]
    rvnf
     & = \Hash(o_{rv},rv)                                 \\[5pt]
    pk_{own}
     & = \Hash(sk_{own})
    \,;\,
    addr=\Hash(pk_{own})
  \end{aligned}
\]

\subsection*{POC Process}

\[
  \begin{aligned}
    \mathbf{s}_{in}
     & := \sum_{i=1}^{m}\mathbf{s}_{in,i} \\[4pt]
    \mathbf{I}
     & := \mathbf{s}_{in}+\mathbf{P}      \\[4pt]
    \mathbf{I}\mathbf{A}
     & := \mathbf{S}_{out}
    =
    \begin{bmatrix}
      \mathbf{s}_{out,1} & \cdots & \mathbf{s}_{out,n}
    \end{bmatrix}
  \end{aligned}
\]

\section{Canonical 공급망 시나리오}

\subsection{참여자와 등록 물품}

\begin{center}
  \begin{tabularx}{\textwidth}{@{}lX@{}}
    \toprule
    참여자               & 역할                                                              \\
    \midrule
    Aluminum Supplier & Aluminum Entry와 Foil Manufacturer 대상 Transfer 수행                \\
    Cathode Supplier  & Cathode Entry와 Cell Manufacturer 대상 Transfer 수행                 \\
    Anode Supplier    & Anode Entry와 Cell Manufacturer 대상 Transfer 수행                   \\
    Foil Manufacturer & 두 제조 공장에서 Foil Process를 수행하고 Foil을 Cell Manufacturer에게 Transfer \\
    Cell Manufacturer & 두 Foil을 Merge하고 Cathode, Anode와 Foil을 사용해 Cell Process 수행       \\
    \bottomrule
  \end{tabularx}
\end{center}

State 값은 $10^9$으로 encode되어 저장된다. 아래 표의 State는 사람이 읽는 단위로
표시한다. 동일한 Product의 개별 note는 Product 이름에 별도 접미사를 붙이지 않고
서로 다른 commitment와 opening으로 구분한다.

\begin{center}
  \begin{tabular}{@{}l@{\hspace{10mm}}llll@{}}
    \toprule
    Entry Product & $q$ & total kg & carbon kgCO2e & recycled kg \\
    \midrule
    Aluminum      & 10  & 20       & 160           & 10          \\
    Aluminum      & 5   & 10       & 150           & 0           \\
    Cathode       & 4   & 200      & 400           & 40          \\
    Anode         & 3   & 100      & 150           & 20          \\
    \bottomrule
  \end{tabular}
\end{center}

\subsection{Event 흐름}

Canonical scenario는 다음 공급망 흐름을 25번의 Contract call로 실행한다.

\begin{enumerate}
  \item Aluminum Supplier가 두 Aluminum을 등록한다. 첫 Transfer는 Recall한 뒤 다시
        Transfer하고, 두 Aluminum 모두 Foil Manufacturer가 Proceed한다.
  \item Foil Manufacturer는 각 Aluminum을 Foil과 Foil Scrap으로 Process한다. 생성된 두
        Foil을 Cell Manufacturer에게 Transfer하고, Cell Manufacturer는 Proceed 후 Merge한다.
  \item Cathode와 Anode도 각 Supplier가 Cell Manufacturer에게 Transfer하고, Cell
        Manufacturer가 Proceed한다.
  \item Cell Manufacturer는 Cathode, Anode와 Merge된 Foil을 Cell, Cell Scrap과 Waste로
        Process한다. Cell을 Split한 뒤 첫 Cell output과 Waste를 Exit한다.
\end{enumerate}

각 Contract call의 protocol object 전이와 측정값은 5장에서 확인한다.

\clearpage
\section{Circuit scaling 결과}

\subsection{Canonical verifier의 circuit 크기}

25-call scenario가 실제로 사용하는 Relation의 constraint 수는 다음과 같다.

\begin{center}
  \begin{tabular}{@{}p{0.32\textwidth}l@{}}
    \toprule
    Relation & Constraints \\
    \midrule
    Entry    & 2,621       \\
    Transfer & 28,843      \\
    Proceed  & 15,855      \\
    Recall   & 15,855      \\
    Merge    & 28,298      \\
    Split    & 27,938      \\
    Process  & 62,833      \\
    Exit     & 13,136      \\
    \bottomrule
  \end{tabular}
\end{center}

Canonical Process 하나가 Foil 1-to-2와 Cell 3-to-3을 모두 처리한다. 이 circuit은 public
$m,n$에서 active slot을 만들고 inactive input/output과 $\mathbf{A}$의 inactive row를 zero로
강제한다.

\subsection{State 길이에 따른 constraint 증가}

아래 표는 State coordinate를 0개에서 3개까지 늘렸을 때의 constraint 수다.

\begin{center}
  \begin{tabular}{@{}llllll@{}}
    \toprule
    \# of State & \multicolumn{1}{l}{0} & \multicolumn{1}{l}{1} & \multicolumn{1}{l}{2} & \multicolumn{1}{l}{3} & Constraints 증가량 \\
    \midrule
    Entry       & 1,334                 & 1,763                 & 2,192                 & 2,621                 & 429             \\
    Transfer    & 15,808                & 20,153                & 24,498                & 28,843                & 4,345           \\
    Proceed     & 14,046                & 14,649                & 15,252                & 15,855                & 603             \\
    Recall      & 14,046                & 14,649                & 15,252                & 15,855                & 603             \\
    Merge       & 25,199                & 26,232                & 27,265                & 28,298                & 1,033           \\
    Split       & 14,903                & 19,248                & 23,593                & 27,938                & 4,345           \\
    Process     & 39,495                & 47,174                & 54,853                & 62,533                & 7,679.3         \\
    Exit        & 12,233                & 12,534                & 12,835                & 13,136                & 301             \\
    \bottomrule
  \end{tabular}
\end{center}

Process의 증가폭이 가장 크다. 3-to-3 Process는 각 State coordinate마다 input/output
commitment binding, aggregate transition, allocation과 두 floor/remainder relation을
추가한다. Transfer와 Split도 proportional floor relation, remainder bound와 exact aggregate
conservation 때문에 증가폭이 크다. Proceed와 Recall은 voucher payload와 output State의
equality를 확인하므로 증가폭이 작다.

3-to-3 Process에서는 State coordinate 하나가 늘어날 때 public delta $P_k$ 하나와
세 output의 allocation $A_{1,k},A_{2,k},A_{3,k}$가 추가된다. 따라서 State coordinate당
public input이 4개 증가한다.

\subsection{Process 입출력 크기에 따른 변화}

State 길이를 3으로 고정하고 circuit의 최대 input/output slot 수를 늘린 결과다. 여기서
Process 입출력 크기 $a$는 compile 시 고정되는 최대 크기
$M_{\max}=N_{\max}=a$를 뜻한다. 예를 들어 3-to-3은 circuit이 input과 output을 각각
최대 3개까지 수용하도록 구성됐다는 의미다.

\begin{center}
  \begin{tabular}{@{}l@{\hspace{28mm}}l@{\hspace{28mm}}l@{}}
    \toprule
    Process 입출력 크기 & Constraints & Prove ms \\
    \midrule
    $M=1,\,N=1$    & 19,309      & 438      \\
    $M=2,\,N=2$    & 40,921      & 841      \\
    $M=3,\,N=3$    & 62,533      & 847      \\
    \bottomrule
  \end{tabular}
\end{center}

최대 입출력 크기가 증가하면 input commitment의 opening과 membership, nullifier, output commitment와
allocation relation이 함께 추가되므로 constraint 수가 거의 선형으로 증가한다. Prove 시간은
단일 smoke run의 참고값이므로 constraint 증가에 정확히 비례하거나 대표 성능으로 해석하지 않는다.
반복 측정한 시간 지표는 4.2의 clean Anvil 5회 E2E 결과를 사용한다.

\section{On-chain gas 결과}

\subsection{Canonical Contract 배포 gas}

각 Contract를 독립 Anvil deployment transaction으로 제출하고 receipt의
\texttt{gasUsed}를 측정했다. \texttt{ZkDPP}는 별도 Merkle Tree Contract를 배포하지 않고
내부에 $\MT$와 $\rvMT$ Tree struct를 가진다. 따라서 Main Contract 배포값에는
constructor가 두 Tree의 depth-32 zero node를 계산하고 저장하는 비용이 포함된다.

\begin{center}
  \begin{tabular}{@{}ll@{}}
    \toprule
    배포 대상                  & Gas        \\
    \midrule
    Poseidon2BLS12381      & 2,282,769  \\
    Entry verifier         & 1,602,240  \\
    Transfer verifier      & 1,600,848  \\
    Proceed verifier       & 1,601,580  \\
    Recall verifier        & 1,601,532  \\
    Merge verifier         & 1,602,240  \\
    Split verifier         & 1,601,568  \\
    Process verifier       & 1,602,372  \\
    Exit verifier          & 1,601,820  \\
    ZkDPP (MT/rvMT 초기화 포함) & 4,908,046  \\
    \midrule
    컴포넌트 합계                & 20,005,015 \\
    \bottomrule
  \end{tabular}
\end{center}

배포 10건을 각각 독립 transaction으로 보냈으며, ZkDPP 배포값은 constructor의
$\MT/\rvMT$ depth-32 초기화를 포함한다.

\clearpage
\subsection{Event별 Gas 및 E2E 시간}

아래 값은 canonical scenario에서 같은 Event가 여러 번 실행된 경우 각 transaction 결과의
median이다. Process는 canonical 최대 입출력 크기인 3-to-3을 실행한 22번 transaction만
사용한다. 한 번만 실행된 Event도 해당 transaction의 측정값을 그대로 사용한다. Gas는
고정 proof를 제출한 Anvil transaction의 \texttt{gasUsed}, E2E ms는 live proof 생성부터
receipt 수신까지 clean Anvil 5회 측정의 median이다.

\begin{center}
  \begin{tabular}{@{}l@{\hspace{25mm}}l@{\hspace{25mm}}l@{}}
    \toprule
    Event    & Gas       & E2E ms  \\
    \midrule
    Entry    & 1,498,510 & 96.725  \\
    Transfer & 2,673,552 & 476.554 \\
    Proceed  & 1,519,990 & 269.926 \\
    Recall   & 1,521,787 & 284.994 \\
    Merge    & 1,561,954 & 459.582 \\
    Split    & 2,352,584 & 476.843 \\
    Process  & 3,315,365 & 846.296 \\
    Exit     & 401,736   & 258.162 \\
    \bottomrule
  \end{tabular}
\end{center}

같은 Event라도 append index와 storage slot의 zero-to-nonzero 전환 여부에 따라 gas가 달라질
수 있다. Merkle Tree leaf 하나를 append할 때 Poseidon2 compression을 32회 실행하며,
Transfer는 $\MT$와 $\rvMT$를 함께 갱신한다. 개별 transaction의 차이는 5.2에서 확인한다.

\clearpage
\section{Canonical scenario 상세 결과}

\subsection{Transaction별 Event 흐름}

5.1과 5.2의 Index는 동일한 canonical Contract call을 가리킨다.

\begin{longtable}{@{}p{0.07\textwidth}p{0.14\textwidth}p{0.69\textwidth}@{}}
  \toprule
  Index & Event    & Protocol object 전이                                                                     \\
  \midrule
  \endhead
  1     & Entry    & $0\rightarrow cm_{\mathrm{Aluminum},1}$                                                \\
  2     & Entry    & $0\rightarrow cm_{\mathrm{Aluminum},2}$                                                \\
  3     & Entry    & $0\rightarrow cm_{\mathrm{Cathode}}$                                                   \\
  4     & Entry    & $0\rightarrow cm_{\mathrm{Anode}}$                                                     \\
  5     & Transfer & $cm_{\mathrm{Aluminum},1}\rightarrow(rv_{\mathrm{Aluminum},1},cm_{\mathrm{change},1})$ \\
  6     & Recall   & $rv_{\mathrm{Aluminum},1}\rightarrow cm_{\mathrm{return}}$                             \\
  7     & Transfer & $cm_{\mathrm{return}}\rightarrow(rv_{\mathrm{Aluminum},2},cm_{\mathrm{change},2})$     \\
  8     & Proceed  & $rv_{\mathrm{Aluminum},2}\rightarrow cm_{\mathrm{Aluminum},3}$                         \\
  9     & Transfer & $cm_{\mathrm{Aluminum},2}\rightarrow(rv_{\mathrm{Aluminum},3},cm_{\mathrm{change},3})$ \\
  10    & Proceed  & $rv_{\mathrm{Aluminum},3}\rightarrow cm_{\mathrm{Aluminum},4}$                         \\
  11    & Process  & $cm_{\mathrm{Aluminum},3}\rightarrow(cm_{\mathrm{Foil},1},cm_{\mathrm{FoilScrap},1})$  \\
  12    & Process  & $cm_{\mathrm{Aluminum},4}\rightarrow(cm_{\mathrm{Foil},2},cm_{\mathrm{FoilScrap},2})$  \\
  13    & Transfer & $cm_{\mathrm{Foil},1}\rightarrow(rv_{\mathrm{Foil},1},cm_{\mathrm{change},4})$         \\
  14    & Proceed  & $rv_{\mathrm{Foil},1}\rightarrow cm'_{\mathrm{Foil},1}$                                \\
  15    & Transfer & $cm_{\mathrm{Foil},2}\rightarrow(rv_{\mathrm{Foil},2},cm_{\mathrm{change},5})$         \\
  16    & Proceed  & $rv_{\mathrm{Foil},2}\rightarrow cm'_{\mathrm{Foil},2}$                                \\
  17    & Merge    & $(cm'_{\mathrm{Foil},1},cm'_{\mathrm{Foil},2})\rightarrow cm_{\mathrm{Foil}}$          \\
  18    & Transfer & $cm_{\mathrm{Cathode}}\rightarrow(rv_{\mathrm{Cathode}},cm_{\mathrm{change},6})$       \\
  19    & Proceed  & $rv_{\mathrm{Cathode}}\rightarrow cm'_{\mathrm{Cathode}}$                              \\
  20    & Transfer & $cm_{\mathrm{Anode}}\rightarrow(rv_{\mathrm{Anode}},cm_{\mathrm{change},7})$           \\
  21    & Proceed  & $rv_{\mathrm{Anode}}\rightarrow cm'_{\mathrm{Anode}}$                                  \\
  22    & Process  & $(cm'_{\mathrm{Cathode}},cm'_{\mathrm{Anode}},cm_{\mathrm{Foil}})
                       \rightarrow(cm_{\mathrm{Cell}},cm_{\mathrm{CellScrap}},cm_{\mathrm{Waste}})$         \\
  23    & Split    & $cm_{\mathrm{Cell}}\rightarrow(cm_{\mathrm{Cell},1},cm_{\mathrm{Cell},2})$             \\
  24    & Exit     & $cm_{\mathrm{Cell},1}\rightarrow0$                                                     \\
  25    & Exit     & $cm_{\mathrm{Waste}}\rightarrow0$                                                      \\
  \bottomrule
\end{longtable}

\subsection{Transaction별 Gas 및 E2E 시간}

Gas는 고정 proof를 사용한 Anvil transaction의 \texttt{gasUsed}다. E2E ms는 live proof를
생성한 clean Anvil 5회의 median이며, 마지막 열은 동일한 5회의 min--max다.

  {\small
    \begin{longtable}{@{}p{0.08\textwidth}p{0.17\textwidth}p{0.19\textwidth}p{0.18\textwidth}p{0.26\textwidth}@{}}
      \toprule
      Index & Event    & Gas       & E2E ms  & E2E min--max ms  \\
      \midrule
      \endhead
      1     & Entry    & 2,056,352 & 97.902  & 93.211--108.453  \\
      2     & Entry    & 1,489,978 & 98.459  & 94.615--102.190  \\
      3     & Entry    & 1,507,042 & 95.548  & 93.475--96.485   \\
      4     & Entry    & 1,487,820 & 95.008  & 93.869--104.164  \\
      5     & Transfer & 3,259,160 & 478.788 & 466.465--499.958 \\
      6     & Recall   & 1,521,787 & 284.994 & 263.890--316.189 \\
      7     & Transfer & 2,673,552 & 477.061 & 473.265--588.660 \\
      8     & Proceed  & 1,517,874 & 268.193 & 260.188--307.967 \\
      9     & Transfer & 2,726,962 & 480.590 & 455.507--493.482 \\
      10    & Proceed  & 1,520,020 & 281.241 & 263.464--298.725 \\
      11    & Process  & 2,385,979 & 841.452 & 824.666--850.305 \\
      12    & Process  & 2,403,091 & 839.126 & 833.020--883.232 \\
      13    & Transfer & 2,669,344 & 472.743 & 458.948--529.966 \\
      14    & Proceed  & 1,515,764 & 269.471 & 260.251--272.064 \\
      15    & Transfer & 2,761,198 & 469.997 & 465.704--492.411 \\
      16    & Proceed  & 1,519,960 & 270.381 & 267.622--276.066 \\
      17    & Merge    & 1,561,954 & 459.582 & 454.967--522.587 \\
      18    & Transfer & 2,652,244 & 476.554 & 470.200--484.047 \\
      19    & Proceed  & 1,554,172 & 266.827 & 258.736--288.917 \\
      20    & Transfer & 2,669,320 & 476.230 & 461.333--503.480 \\
      21    & Proceed  & 1,535,022 & 270.534 & 261.469--281.826 \\
      22    & Process  & 3,315,365 & 846.296 & 827.337--854.559 \\
      23    & Split    & 2,352,584 & 476.843 & 463.276--605.012 \\
      24    & Exit     & 401,742   & 261.554 & 252.118--303.593 \\
      25    & Exit     & 401,730   & 254.769 & 251.522--271.306 \\
      \bottomrule
    \end{longtable}
  }

25개 Event transaction의 Gas 합계는 49,460,016이다. E2E 측정 경계는 Event witness
생성 직전부터 live proof 생성, transaction 제출과 성공 receipt 수신까지다.

\section{Open Problems}

\subsection{Provenance}

현재 POC는 $\MT$와 $\rvMT$에 leaf를 append할 때 leaf, index와 new root를 event로
발행한다. 이 event는 Client의 Merkle path 조회를 위한 것이며, 소비된 input $cm$과
생성된 output $cm$ 사이의 provenance edge를 직접 선언하지 않는다.

미해결 질문은 input/output $cm$을 연결하는 provenance event를 발행할 것인지다.

\begin{itemize}
  \item Event를 발행하면 commitment 간 linking, product graph와 $cm$ spent time이
        공개된다.
  \item 직접 downstream뿐 아니라 edge를 재귀적으로 따라간 간접 downstream도 추적할 수
        있다.
  \item 현재 calldata에도 input/output $cm$이 포함되므로, event를 발행하지 않아도 동일한
        linking이 가능한지 확인해야 한다.
\end{itemize}

DocumentInfo와 State 값을 숨기는 것과 product graph 및 timing을 숨기는 것은 별개의
privacy 문제다. Provenance를 제공할 범위와 논문에서 주장할 privacy 경계를 함께 결정해야
한다.

\subsection{Recipe 기반 Process}

현재 Process는 input State를 합산하고 public process delta $\mathbf{P}$와
StateAllocationMatrix $\mathbf{A}$를 적용했을 때 output State가 선언된 relation과
일치하는지만 증명한다. Output quantity는 range와 commitment binding만 검증하며,
input/output Material Type과 Quantity transition은 강제하지 않는다. 따라서 호출자는
circuit 식을 만족하는 범위에서 $\mathbf{P}$와 $\mathbf{A}$를 임의로 선택할 수 있다.

Recipe 기반 Process는 다음 조건을 함께 강제해야 한다.

\begin{itemize}
  \item 허용되는 input/output Material Type
  \item input quantity에서 output quantity로의 transition rule
  \item 허용되는 State delta
  \item output별 State allocation policy
\end{itemize}

Material Type을 검증하려면 MaterialClassId를 note에 binding하거나 DocumentHash와
Material Type을 연결하는 별도 proof가 필요하다. 이때 Process를 거친 물품의
input/output Material 관계가 공개될 수 있고, 유효한 Recipe를 등록하고 관리할
신뢰 주체도 필요하다. Recipe privacy와 Recipe registry의 trust model을 함께 해결해야
한다.

\section{결과 해석}

\begin{enumerate}
  \item \textbf{State transition consistency는 실제 8-Event chain에서 동작했다.}
        Merge/Split 합 보존, Transfer 비례 배분, Proceed/Recall의 voucher 단일 소비,
        Foil 1-to-2 및 Cell 3-to-3 Process의 public delta/allocation과 Exit 소비가 같은
        scenario에서 검증됐다.
  \item \textbf{비례 배분이 가장 큰 circuit 증가 요인이다.}
        Transfer와 Split은 State coordinate당 4,345 constraints, 3-to-3 Process는 약
        7,679 constraints가 증가했다.
  \item \textbf{전체 Event gas는 storage update가 크게 좌우한다.}
        Exit은 append가 없어 약 402K gas인 반면, 두 Tree를 갱신하는 Transfer는
        2.65--3.26M gas, output 두 개를 append하는 Foil Process는 약 2.39--2.40M gas,
        output 세 개를 append하는 Cell Process는 약 3.32M gas를 사용했다.
  \item \textbf{Prover 시간과 Event gas는 같은 비율로 증가하지 않는다.}
        Circuit 규모 증가는 주로 off-chain proof 생성 시간에 반영된다. On-chain Event
        gas는 proof verification뿐 아니라 Tree append와 storage write의 영향을 크게 받는다.
  \item \textbf{E2E 시간은 proof 생성이 지배한다.}
        Foil Process 두 건의 E2E median은 841.452ms와 839.126ms였고, Cell Process는
        전체 E2E 846.296ms였다.
        E2E는 clean Anvil 5회 median과 min/max를 함께 기록했다.
\end{enumerate}

\section{재현 방법}

\begin{verbatim}
cd poc-v2
make test-go
make setup
make test-contract
make benchmark
make benchmark-anvil-gas ANVIL_PORT=18545
make benchmark-anvil-e2e RUNS=5 ANVIL_PORT=18545
make benchmark-anvil RUNS=5 ANVIL_PORT=18545
\end{verbatim}

\texttt{make benchmark}는 34개 profile의 Go compile/setup/prove/verify 결과와 Solidity
verifier fixture를 생성한다.
\texttt{make benchmark-anvil-gas}는 고정 proof를 사용해 실제 Anvil gas를 생성한다.
\texttt{make benchmark-anvil-e2e}는 각 Event proof를 실행 중 생성하며 clean chain을 다섯 번
반복한다. Host port 8545가 이미 사용 중이면 위 예시처럼 \texttt{ANVIL\_PORT}를 변경한다.

Canonical input을 변경할 때는 \path{testdata/scenario/canonical.json}만 수정하고
\texttt{make generate-reference}로 derived fixture를 다시 생성한다. Event별 숫자를 별도
fixture에 수작업으로 중복 정의하지 않는다.

\end{document}
